\begin{center}
    {\LARGE \textbf{2007物理化学}} \\[2ex]
    {\large 时量180分钟 \quad 满分90分}
\end{center}

\vspace{0.5cm}
\noindent\textbf{注意事项:}
\begin{enumerate}[label=\arabic*., parsep=0pt, itemsep=0pt]
    \item 答题前,考生先将自己的姓名、学号、考场号、座位号填写在答题卡上,将条形码准确粘贴在考生信息条形码粘贴区。
    \item 考试结束后,将本试卷与答题纸一并收回。
    \item 回答各个问题时,将答案写在答题纸上,写在本试卷上无效。
\end{enumerate}

\vspace{0.5cm}
\noindent\textbf{1. (15 pts.)}

取$0^\circ\mathrm{C}$,$3p^\ominus$的$\ce{O_{2}(g)}$ $10\ \mathrm{dm^3}$,绝热膨胀到压力为$p^\ominus$,分别计算下列两种过程的$\Delta G$:
\begin{enumerate}[label=(\arabic*), itemsep=1ex]
    \item 绝热可逆膨胀；
    \item 将外压力骤减至$p^\ominus$,气体反抗外压力进行绝热不可逆膨胀。
\end{enumerate}
假定$\ce{O_{2}(g)}$为理想气体,其摩尔定容热容$C_{V,\mathrm{m}}=\dfrac{5}{2}R$。已知氧气的摩尔标准熵$S_\mathrm{m}^\ominus(298\ \mathrm{K})=205.0\ \mathrm{J\cdot K^{-1}\cdot mol^{-1}}$。

\vspace{0.5cm}
\noindent\textbf{2. (15 pts.)}

$300\ \mathrm{K}$时,液态$\ce{A}$的蒸气压为$37338\ \mathrm{Pa}$,液态$\ce{B}$的蒸气压为$22656\ \mathrm{Pa}$。当$2\ \mathrm{mol}$ $\ce{A}$和$2\ \mathrm{mol}$ $\ce{B}$混合后,液面上蒸气压为$50663\ \mathrm{Pa}$,在蒸气中$\ce{A}$的摩尔分数为$0.60$,假定蒸气为理想气体,以纯液态$\ce{A}$和$\ce{B}$为标准态。
\begin{enumerate}[label=(\arabic*), itemsep=1ex]
    \item 求溶液中$\ce{A}$和$\ce{B}$的活度；
    \item 求溶液中$\ce{A}$和$\ce{B}$的活度因子；
    \item 求$\ce{A}$和$\ce{B}$的混合Gibbs自由能$\Delta_\mathrm{mix}G$。
\end{enumerate}

\vspace{0.5cm}
\noindent\textbf{3. (15 pts.)}

已知萘的熔点为$353.2\ \mathrm{K}$,在此温度时的摩尔熔化热$\Delta_\mathrm{fus}H_\mathrm{m}=18.995\ \mathrm{kJ\cdot mol^{-1}}$。计算萘在$313.2\ \mathrm{K}$时与某溶剂形成理想液体混合物时的溶解度(以摩尔分数表示)。

\vspace{0.5cm}
\noindent\textbf{4. (15 pts.)}

在真空的容器中放入固态$\ce{NH_{4}HS}$,于$25^\circ\mathrm{C}$下分解为$\ce{NH_{3}}$和$\ce{H_{2}S}$,平衡时容器内的压力为$6.665\times10^4\ \mathrm{Pa}$。
\begin{enumerate}[label=(\arabic*), itemsep=1ex]
    \item 若放入$\ce{NH_{4}HS}$时容器中已有$3.998\times10^4\ \mathrm{Pa}$的$\ce{H_{2}S}$,求平衡时容器中的压力；
    \item 若容器中原有$6.665\times10^3\ \mathrm{Pa}$的$\ce{NH_{3}}$,问需加多大压力的$\ce{H_{2}S}$才能开始形成$\ce{NH_{4}HS}$固体。
\end{enumerate}

\vspace{0.5cm}
\noindent\textbf{5. (15 pts.)}

阿司匹林的热分解作用是一级反应,$458\ \mathrm{K}$和$510\ \mathrm{K}$时的速率常数分别为$2.2\times10^{-5}$和$3.07\times10^{-3}\ \mathrm{min^{-1}}$,试求反应的实验活化能$E_\mathrm{a}$以及在平均温度时的活化焓和活化熵。

\vspace{0.5cm}
\noindent\textbf{6. (15 pts.)}

$298\ \mathrm{K}$时,有下列电池:
\[
\ce{Pt,Cl_{2}(p^\ominus)}|\ce{HCl(0.1\ mol\cdot kg^{-1})}|\ce{AgCl(s)}|\ce{Ag(s)}
\]
试求:
\begin{enumerate}[label=(\arabic*), itemsep=1ex]
    \item 电池的电动势；
    \item 电动势温度系数和有$1\ \mathrm{mol}$电子电荷量可逆输出时的热效应；
    \item $\ce{AgCl(s)}$的分解压力。
\end{enumerate}
已知$\Delta_\mathrm{f}H_\mathrm{m}^\ominus(\ce{AgCl})=-1.2703\times10^5\ \mathrm{J\cdot mol^{-1}}$,$\ce{Ag(s)}$,$\ce{AgCl(s)}$和$\ce{Cl_{2}(g)}$的规定熵值$S_\mathrm{m}^\ominus$分别为$42.70\ \mathrm{J\cdot K^{-1}\cdot mol^{-1}}$,$96.11\ \mathrm{J\cdot K^{-1}\cdot mol^{-1}}$和$243.87\ \mathrm{J\cdot K^{-1}\cdot mol^{-1}}$。

